\section{Background}
\label{sec:background}

In the C++ language, a program that calls \texttt{new} first asks an allocator for a block of free memory.
The allocator then finds a free block large enough, records that the block is in
use, and returns a pointer to it.
\texttt{delete} will then return the block back so it can be used again.
The base allocator supplied with the operating system serves every program on the
machine, so it is built for the average case rather than for any one memory access
pattern.
Performance costs related to memory allocation can come from searching for a block that fits, from locking shared
structures so that threads do not corrupt them, from the gaps left between
blocks as memory is reused, and from the work the operating system does to
supply memory as a program grows.
Searching longer for a tighter fit wastes less space but takes more time, and
holding more free memory in reserve avoids requests to the operating system at
the cost of a larger footprint.
When the memory access pattern of a program is known, opting for a custom memory allocator
can improve performance by tuning the structure and algorithm of allocation to fit it. 
